% TOP_PROBLEM: {"id":30007747,"problem_number":"LOCAL-30007747","title":"Gentrification and incumbent welfare","category_id":21,"category":{"id":21,"name":"economics","display_name":"Economics","description":"Open economic theory statements, published research agendas, and proposed scoped specifications; question provenance and historical priority are qualified per record.","slug":"economics","order_index":21,"created_at":"2026-10-08T00:00:00Z"},"status":"research_question","external_url":null,"legacy_ids":[],"published":false,"statement":"Estimate ten-year welfare and involuntary displacement for all baseline residents of neighborhoods receiving a defined revitalization package, tracking those who leave.","economics":{"original_id":236,"field":"Urban, housing and spatial economics","kind":"empirical","formalization_basis":"adapted_research_question","source_status":"explicit_open","priority_status":"earliest_found","priority_note":"Earliest explicit statement located in this audit; no exhaustive priority search or first-ever claim. The mathematical specification below is adapted, not quoted from the source.","earliest_verified_reference":{"authors":["Eric Chyn","Lawrence F. Katz"],"title":"Neighborhoods Matter: Assessing the Evidence for Place Effects","year":2021,"url":"https://lkatz.scholars.harvard.edu/sites/g/files/omnuum5961/files/lkatz/files/chyn_katz_jep_2021_all.pdf","doi":"10.1257/jep.35.4.197","locator":"Conclusion, printed pp. 216–217; introduction p. 200","evidence_summary":"The authors explicitly identify mechanism weights, general-equilibrium scaling of mobility programs, and effects of revitalization on preexisting residents as open issues."},"current_reference":{"authors":["Eric Chyn","Lawrence F. Katz"],"title":"Neighborhoods Matter: Assessing the Evidence for Place Effects","year":2021,"url":"https://lkatz.scholars.harvard.edu/sites/g/files/omnuum5961/files/lkatz/files/chyn_katz_jep_2021_all.pdf","doi":"10.1257/jep.35.4.197","locator":"Conclusion, printed pp. 216–217; introduction p. 200","evidence_summary":"The authors explicitly identify mechanism weights, general-equilibrium scaling of mobility programs, and effects of revitalization on preexisting residents as open issues."},"formalization_note":"The source explicitly identifies the underlying gap or agenda. Population, horizon, interventions, estimands, and auxiliary assumptions are proposed operational choices, not a canonical source theorem. Partial later answers are compatible with this scoped question; the citation alone does not prove absence of every subsequent solution. The 2021 Chyn–Katz review supplies an earlier verified agenda than the original catalogue’s 2023 or 2026 supporting citation."},"statusQualification":"Published question or research agenda verified at the cited locator. The mathematical specification is adapted for this catalogue and includes additional assumptions; source publication does not establish first-ever priority or certify this exact model remains unsolved. The source explicitly identifies the underlying gap or agenda. Population, horizon, interventions, estimands, and auxiliary assumptions are proposed operational choices, not a canonical source theorem. Partial later answers are compatible with this scoped question; the citation alone does not prove absence of every subsequent solution. The 2021 Chyn–Katz review supplies an earlier verified agenda than the original catalogue’s 2023 or 2026 supporting citation.","statusReviewedAt":"2026-10-08"}
% FIRST_POSTED: 2026-10-08
% ENTRY_KIND: statement_only
% !TeX program = lualatex
\documentclass[11pt]{article}
\usepackage{fontspec}
\usepackage[a4paper,margin=25mm]{geometry}
\usepackage{amsmath,amssymb,mathtools}
\usepackage{xurl}
\usepackage[hidelinks]{hyperref}
\newcommand{\sref}[2]{\hyperlink{source:#1:#2}{[S#2]}}
\setlength{\emergencystretch}{3em}
\begin{document}
% problemId:30007747
\section*{Gentrification and incumbent welfare}\label{30007747}
\subsection{Definitions and mathematical statement}
Consider households living in a prespecified set of low-income urban neighborhoods immediately before a revitalization program is announced. Let \(a=1\) denote the defined investment package and \(a=0\) its absence. For baseline household \(i\), let \(V_i(a)\) be ten-year monetary-equivalent welfare combining consumption, housing quality, travel, safety, and valued neighborhood ties under the realized residential path. Define displacement \(D_i(a)\in\{0,1\}\) as a move caused by eviction or an unaffordable housing cost increase. The targets are
\[\tau_V=E[V_i(1)-V_i(0)],\qquad\tau_D=E[D_i(1)-D_i(0)].\]
Follow every baseline household after moving; repeated cross-sectional neighborhood averages do not identify these incumbent outcomes. Specify credible program variation, announcement effects, and counterfactual trends, and measure nearby spillovers and landlord capital gains separately. Distinguish voluntary moves from displacement using a prespecified classification with sensitivity analysis. Monetizing social ties requires explicit preference assumptions or bounds, rather than treating every move as a welfare loss. The Chyn–Katz review explicitly identifies benefits to preexisting residents of revitalization as unclear. This adapted question focuses on one intervention and horizon; it does not claim that all forms of gentrification or their effects are unknown.

\par\textbf{Formulation and scope.} The source explicitly identifies the underlying gap or agenda. Population, horizon, interventions, estimands, and auxiliary assumptions are proposed operational choices, not a canonical source theorem. Partial later answers are compatible with this scoped question; the citation alone does not prove absence of every subsequent solution. The 2021 Chyn–Katz review supplies an earlier verified agenda than the original catalogue’s 2023 or 2026 supporting citation.

\par\textbf{Open-question provenance.} An explicit question or research agenda was verified at the source locator. This does not by itself certify that every later version remains unresolved \sref{30007747}{1}.

\par\textbf{Historical priority.} Earliest explicit statement located in this audit; no exhaustive priority search or first-ever claim. The mathematical specification below is adapted, not quoted from the source.

\subsection{Short English statement}
Estimate ten-year welfare and involuntary displacement for all baseline residents of neighborhoods receiving a defined revitalization package, tracking those who leave.

\subsection{Sources}
\begin{itemize}\raggedright
\item[\textnormal{[S1]}]\hypertarget{source:30007747:1}{} Eric Chyn, Lawrence F. Katz. 2021. Neighborhoods Matter: Assessing the Evidence for Place Effects. Conclusion, printed pp. 216–217; introduction p. 200.
\url{https://lkatz.scholars.harvard.edu/sites/g/files/omnuum5961/files/lkatz/files/chyn_katz_jep_2021_all.pdf}.
DOI: 10.1257/jep.35.4.197.
{\small\textit{Earliest verified question/agenda reference; historical priority qualified below. The authors explicitly identify mechanism weights, general-equilibrium scaling of mobility programs, and effects of revitalization on preexisting residents as open issues.}}
\end{itemize}
\end{document}
